\documentclass[11pt,a4paper]{article}
\usepackage[utf8]{inputenc}
\usepackage[T1]{fontenc}
\usepackage[spanish,es-minimal,es-noshorthands]{babel}
\AtBeginDocument{\renewcommand{\sin}{\operatorname{sen}}\renewcommand{\arcsin}{\operatorname{arcsen}}}
\usepackage{lmodern}
\usepackage[sfdefault,lining]{FiraSans}
\usepackage[a4paper,top=2.3cm,bottom=2.2cm,left=1.8cm,right=1.8cm,headheight=1.3cm,headsep=0.4cm,footskip=1cm]{geometry}
\usepackage{amsmath,amssymb}
\usepackage{icomma}
\usepackage{graphicx}
\usepackage[table]{xcolor}
\usepackage{tikz}
\usepackage{fontawesome5}
\usepackage{ccicons}
\usepackage{fancyhdr}
\usepackage{lastpage}
\usepackage{tasks}
\usepackage{needspace}
\usepackage[most]{tcolorbox}
\usepackage{array}
\usepackage[hidelinks]{hyperref}
\AtBeginDocument{\hypersetup{pdftitle={Ficha de ejercicios de la unidad 1: Números reales. 1º Bach.},pdfsubject={Matemáticas CCSS 1º Bach.},pdfkeywords={números reales, fracción generatriz, intervalos, semirrectas, valor absoluto, unión e intersección, radicales, forma exponencial, índice común, racionalización, logaritmos}}}

\definecolor{navy}{HTML}{1B2748}
\definecolor{accent}{HTML}{E8603A}
\definecolor{teal}{HTML}{1F9E90}
\definecolor{soft}{HTML}{EEF1F8}
\definecolor{muted}{HTML}{5D6685}

\setlength{\parindent}{0pt}
\setlength{\parskip}{3pt}
\settasks{label=\alph*), label-format=\bfseries\color{navy}, label-width=1.3em, item-indent=1.8em, column-sep=0.6em, after-item-skip=5pt, before-skip=3pt, after-skip=4pt}

% ---------- capçalera i peu ----------
\newcommand{\llicencia}{\href{https://creativecommons.org/licenses/by-nc-sa/4.0/deed.es}{\ccbyncsa}\ \ Licencia Creative Commons Atribución-NoComercial-CompartirIgual 4.0 Internacional (CC BY-NC-SA 4.0)}
\pagestyle{fancy}
\fancyhf{}
\fancyhead[L]{\raisebox{-0.35\height}{\includegraphics[height=0.95cm]{logo-gva.png}}}
\fancyhead[C]{\small\color{muted}Ficha de ejercicios · Unidad 1: Números reales · 1\textsuperscript{o} Bach.}
\fancyhead[R]{\raisebox{-0.35\height}{\includegraphics[height=0.95cm]{logo-vermellar.png}}}
\fancyfoot[L]{\scriptsize\color{muted}\llicencia}
\fancyfoot[R]{\small\color{muted}\thepage}
\renewcommand{\headrulewidth}{0.4pt}
\renewcommand{\headrule}{\color{muted!40}\hrule height \headrulewidth}
\fancypagestyle{primera}{\fancyhf{}\renewcommand{\headrulewidth}{0pt}
  \fancyfoot[L]{\scriptsize\color{muted}\llicencia}
  \fancyfoot[R]{\small\color{muted}\thepage}}

% ---------- nivells i exercicis ----------
\newcommand{\estrelles}[1]{{\color{accent}\ifcase#1\or\faStar\faStar[regular]\faStar[regular]\or\faStar\faStar\faStar[regular]\or\faStar\faStar\faStar\fi}}
\newcounter{exer}
\newcommand{\ex}[1]{\par\needspace{4\baselineskip}\medskip\refstepcounter{exer}\noindent{\color{navy}\bfseries\theexer.}\ \ }

\newcommand{\seccio}[2]{\par\needspace{16\baselineskip}\bigskip
  \begin{tcolorbox}[colback=navy,colframe=navy,coltext=white,arc=2mm,boxrule=0pt,left=3mm,right=3mm,top=1.2mm,bottom=1.2mm]
  {\large\bfseries #1\quad #2}
  \end{tcolorbox}\vspace{-2pt}}

\newcommand{\exf}[1]{\par\medskip\refstepcounter{exer}\noindent{\color{navy}\bfseries\theexer.}\ \ }
\newcolumntype{C}[1]{>{\centering\arraybackslash}p{#1}}
% quadre de recordatori (només a la fitxa d'exercicis)
\newcommand{\recorda}[1]{\par\smallskip\begin{tcolorbox}[colback=accent!7,colframe=accent!55,arc=2mm,boxrule=0.5pt,left=2.5mm,right=2.5mm,top=1mm,bottom=1mm]\small{\color{accent}\faLightbulb[regular]}\ \textbf{Recuerda:}\ #1\end{tcolorbox}}
\newcommand{\res}[1]{#1}\newcommand{\rest}[1]{#1}

\def\fmneed{11}% que el títol d'una secció no quede sol al final del full
% =====================================================================
% Capçalera FentMates de les fitxes: un color per a cada curs.
% Els colors són els de la pàgina inicial (LEVELS a assets/js/materials.js);
% el logo ve de _eines/logo/pdf/ (python3 _eines/logo/genera.py).
%
% Ús, al final del preàmbul de la unitat:
%   % =====================================================================
% Capçalera FentMates de les fitxes: un color per a cada curs.
% Els colors són els de la pàgina inicial (LEVELS a assets/js/materials.js);
% el logo ve de _eines/logo/pdf/ (python3 _eines/logo/genera.py).
%
% Ús, al final del preàmbul de la unitat:
%   % =====================================================================
% Capçalera FentMates de les fitxes: un color per a cada curs.
% Els colors són els de la pàgina inicial (LEVELS a assets/js/materials.js);
% el logo ve de _eines/logo/pdf/ (python3 _eines/logo/genera.py).
%
% Ús, al final del preàmbul de la unitat:
%   \input{../comu/fentmates.tex}
%   \fmcurs{1eso}
% i, al començament del document:
%   \capcalera{Matemàtiques}{1r ESO}{Fitxa d'exercicis · Unitat 2}{Divisibilitat}
% Opcional abans de \input: \def\fmneed{14} (espai mínim abans de cada secció).
% =====================================================================
\makeatletter
\def\fm@logodir{../../logo/pdf}
\providecommand{\fmneed}{10}
\def\fm@cursnom{}

% clau del curs, color i color del text damunt del color
\newcommand{\fm@curs}[3]{\@namedef{fm@color@#1}{#2}\@namedef{fm@text@#1}{#3}}
\fm@curs{1eso}{E6B84A}{navy}
\fm@curs{2eso}{3B82F6}{white}
\fm@curs{3eso}{FF7A50}{white}
\fm@curs{4eso-a}{A78BFA}{white}
\fm@curs{4eso-b}{6D28D9}{white}
\fm@curs{1batx-mi}{9B1C1C}{white}
\fm@curs{1batx-ccss}{EF6B6B}{white}
\fm@curs{2batx-mii}{1E6B47}{white}
\fm@curs{2batx-ccss}{2BB3A3}{white}

\newcommand{\fmcurs}[1]{%
  \@ifundefined{fm@color@#1}{\PackageError{fentmates}{Curs desconegut: #1}{}}{}%
  \edef\fm@tmp{\noexpand\definecolor{curs}{HTML}{\@nameuse{fm@color@#1}}%
               \noexpand\colorlet{curstext}{\@nameuse{fm@text@#1}}}\fm@tmp
  \colorlet{cursfosc}{curs!72!black}%
  \def\fm@logo{\fm@logodir/vertical-#1.pdf}%
  \def\fm@marca{\fm@logodir/marca-#1.pdf}}

% ---------- capçalera de la primera pàgina ----------
% #1 matèria (petit), #2 curs, #3 tipus de document i unitat, #4 títol
\newcommand{\capcalera}[4]{%
  \gdef\fm@cursnom{#2}%
  \par\noindent
  \begin{tcolorbox}[colback=curs,colframe=curs,boxrule=0pt,arc=2.2mm,width=3.5cm,height=2.15cm,
      valign=center,left=3mm,right=2mm,top=1mm,bottom=1mm,nobeforeafter,box align=center]
    \color{curstext}{\footnotesize #1}\par\vspace{1pt}{\fontsize{19}{21}\selectfont\bfseries #2\par}%
  \end{tcolorbox}\hspace{5mm}%
  \begin{minipage}[c]{\dimexpr\linewidth-3.5cm-5mm-2.9cm\relax}
    {\color{cursfosc}\small\bfseries\MakeUppercase{#3}}\par\vspace{2pt}
    {\color{navy}\fontsize{20}{23}\selectfont\bfseries #4\par}
  \end{minipage}\hfill
  \begin{minipage}[c]{2.6cm}\raggedleft\includegraphics[width=2.6cm]{\fm@logo}\end{minipage}\par
  \vspace{2pt}}

% ---------- barres de secció del color del curs ----------
\renewcommand{\seccio}[2]{\par\needspace{\fmneed\baselineskip}\bigskip
  \begin{tcolorbox}[colback=curs,colframe=curs,coltext=curstext,arc=2mm,boxrule=0pt,left=3mm,right=3mm,top=1.2mm,bottom=1.2mm]
  {\large\bfseries #1\quad #2}
  \end{tcolorbox}\vspace{-2pt}}

% ---------- estrelles de dificultat del color del curs (a Batxillerat no se'n posen) ----------
\@ifundefined{estrelles}{}{%
  \renewcommand{\estrelles}[1]{{\color{curs}\ifcase#1\or\faStar\faStar[regular]\faStar[regular]\or\faStar\faStar\faStar[regular]\or\faStar\faStar\faStar\fi}}}

% ---------- capçalera i peu de les pàgines ----------
% primera pàgina: logos institucionals al peu
\fancypagestyle{primera}{\fancyhf{}\renewcommand{\headrulewidth}{0pt}
  \fancyfoot[L]{\raisebox{-0.32\height}{\includegraphics[height=0.78cm]{logo-gva.png}}\hspace{5mm}%
    \raisebox{-0.32\height}{\includegraphics[height=0.78cm]{logo-vermellar.png}}}
  \fancyfoot[R]{\scriptsize\color{muted}\href{https://creativecommons.org/licenses/by-nc-sa/4.0/deed.ca}{\ccbyncsa}\ \ CC BY-NC-SA 4.0 · fentmates.org\qquad\small\thepage}}
% la resta: el símbol a l'esquerra i el curs a la dreta, amb el títol al centre com abans
\fancyhead[L]{\raisebox{-0.3\height}{\includegraphics[height=0.62cm]{\fm@marca}}}
\fancyhead[R]{\small\bfseries\color{cursfosc}\fm@cursnom}
\renewcommand{\headrule}{\color{curs!45}\hrule height \headrulewidth}
\makeatother

%   \fmcurs{1eso}
% i, al començament del document:
%   \capcalera{Matemàtiques}{1r ESO}{Fitxa d'exercicis · Unitat 2}{Divisibilitat}
% Opcional abans de \input: \def\fmneed{14} (espai mínim abans de cada secció).
% =====================================================================
\makeatletter
\def\fm@logodir{../../logo/pdf}
\providecommand{\fmneed}{10}
\def\fm@cursnom{}

% clau del curs, color i color del text damunt del color
\newcommand{\fm@curs}[3]{\@namedef{fm@color@#1}{#2}\@namedef{fm@text@#1}{#3}}
\fm@curs{1eso}{E6B84A}{navy}
\fm@curs{2eso}{3B82F6}{white}
\fm@curs{3eso}{FF7A50}{white}
\fm@curs{4eso-a}{A78BFA}{white}
\fm@curs{4eso-b}{6D28D9}{white}
\fm@curs{1batx-mi}{9B1C1C}{white}
\fm@curs{1batx-ccss}{EF6B6B}{white}
\fm@curs{2batx-mii}{1E6B47}{white}
\fm@curs{2batx-ccss}{2BB3A3}{white}

\newcommand{\fmcurs}[1]{%
  \@ifundefined{fm@color@#1}{\PackageError{fentmates}{Curs desconegut: #1}{}}{}%
  \edef\fm@tmp{\noexpand\definecolor{curs}{HTML}{\@nameuse{fm@color@#1}}%
               \noexpand\colorlet{curstext}{\@nameuse{fm@text@#1}}}\fm@tmp
  \colorlet{cursfosc}{curs!72!black}%
  \def\fm@logo{\fm@logodir/vertical-#1.pdf}%
  \def\fm@marca{\fm@logodir/marca-#1.pdf}}

% ---------- capçalera de la primera pàgina ----------
% #1 matèria (petit), #2 curs, #3 tipus de document i unitat, #4 títol
\newcommand{\capcalera}[4]{%
  \gdef\fm@cursnom{#2}%
  \par\noindent
  \begin{tcolorbox}[colback=curs,colframe=curs,boxrule=0pt,arc=2.2mm,width=3.5cm,height=2.15cm,
      valign=center,left=3mm,right=2mm,top=1mm,bottom=1mm,nobeforeafter,box align=center]
    \color{curstext}{\footnotesize #1}\par\vspace{1pt}{\fontsize{19}{21}\selectfont\bfseries #2\par}%
  \end{tcolorbox}\hspace{5mm}%
  \begin{minipage}[c]{\dimexpr\linewidth-3.5cm-5mm-2.9cm\relax}
    {\color{cursfosc}\small\bfseries\MakeUppercase{#3}}\par\vspace{2pt}
    {\color{navy}\fontsize{20}{23}\selectfont\bfseries #4\par}
  \end{minipage}\hfill
  \begin{minipage}[c]{2.6cm}\raggedleft\includegraphics[width=2.6cm]{\fm@logo}\end{minipage}\par
  \vspace{2pt}}

% ---------- barres de secció del color del curs ----------
\renewcommand{\seccio}[2]{\par\needspace{\fmneed\baselineskip}\bigskip
  \begin{tcolorbox}[colback=curs,colframe=curs,coltext=curstext,arc=2mm,boxrule=0pt,left=3mm,right=3mm,top=1.2mm,bottom=1.2mm]
  {\large\bfseries #1\quad #2}
  \end{tcolorbox}\vspace{-2pt}}

% ---------- estrelles de dificultat del color del curs (a Batxillerat no se'n posen) ----------
\@ifundefined{estrelles}{}{%
  \renewcommand{\estrelles}[1]{{\color{curs}\ifcase#1\or\faStar\faStar[regular]\faStar[regular]\or\faStar\faStar\faStar[regular]\or\faStar\faStar\faStar\fi}}}

% ---------- capçalera i peu de les pàgines ----------
% primera pàgina: logos institucionals al peu
\fancypagestyle{primera}{\fancyhf{}\renewcommand{\headrulewidth}{0pt}
  \fancyfoot[L]{\raisebox{-0.32\height}{\includegraphics[height=0.78cm]{logo-gva.png}}\hspace{5mm}%
    \raisebox{-0.32\height}{\includegraphics[height=0.78cm]{logo-vermellar.png}}}
  \fancyfoot[R]{\scriptsize\color{muted}\href{https://creativecommons.org/licenses/by-nc-sa/4.0/deed.ca}{\ccbyncsa}\ \ CC BY-NC-SA 4.0 · fentmates.org\qquad\small\thepage}}
% la resta: el símbol a l'esquerra i el curs a la dreta, amb el títol al centre com abans
\fancyhead[L]{\raisebox{-0.3\height}{\includegraphics[height=0.62cm]{\fm@marca}}}
\fancyhead[R]{\small\bfseries\color{cursfosc}\fm@cursnom}
\renewcommand{\headrule}{\color{curs!45}\hrule height \headrulewidth}
\makeatother

%   \fmcurs{1eso}
% i, al començament del document:
%   \capcalera{Matemàtiques}{1r ESO}{Fitxa d'exercicis · Unitat 2}{Divisibilitat}
% Opcional abans de \input: \def\fmneed{14} (espai mínim abans de cada secció).
% =====================================================================
\makeatletter
\def\fm@logodir{../../logo/pdf}
\providecommand{\fmneed}{10}
\def\fm@cursnom{}

% clau del curs, color i color del text damunt del color
\newcommand{\fm@curs}[3]{\@namedef{fm@color@#1}{#2}\@namedef{fm@text@#1}{#3}}
\fm@curs{1eso}{E6B84A}{navy}
\fm@curs{2eso}{3B82F6}{white}
\fm@curs{3eso}{FF7A50}{white}
\fm@curs{4eso-a}{A78BFA}{white}
\fm@curs{4eso-b}{6D28D9}{white}
\fm@curs{1batx-mi}{9B1C1C}{white}
\fm@curs{1batx-ccss}{EF6B6B}{white}
\fm@curs{2batx-mii}{1E6B47}{white}
\fm@curs{2batx-ccss}{2BB3A3}{white}

\newcommand{\fmcurs}[1]{%
  \@ifundefined{fm@color@#1}{\PackageError{fentmates}{Curs desconegut: #1}{}}{}%
  \edef\fm@tmp{\noexpand\definecolor{curs}{HTML}{\@nameuse{fm@color@#1}}%
               \noexpand\colorlet{curstext}{\@nameuse{fm@text@#1}}}\fm@tmp
  \colorlet{cursfosc}{curs!72!black}%
  \def\fm@logo{\fm@logodir/vertical-#1.pdf}%
  \def\fm@marca{\fm@logodir/marca-#1.pdf}}

% ---------- capçalera de la primera pàgina ----------
% #1 matèria (petit), #2 curs, #3 tipus de document i unitat, #4 títol
\newcommand{\capcalera}[4]{%
  \gdef\fm@cursnom{#2}%
  \par\noindent
  \begin{tcolorbox}[colback=curs,colframe=curs,boxrule=0pt,arc=2.2mm,width=3.5cm,height=2.15cm,
      valign=center,left=3mm,right=2mm,top=1mm,bottom=1mm,nobeforeafter,box align=center]
    \color{curstext}{\footnotesize #1}\par\vspace{1pt}{\fontsize{19}{21}\selectfont\bfseries #2\par}%
  \end{tcolorbox}\hspace{5mm}%
  \begin{minipage}[c]{\dimexpr\linewidth-3.5cm-5mm-2.9cm\relax}
    {\color{cursfosc}\small\bfseries\MakeUppercase{#3}}\par\vspace{2pt}
    {\color{navy}\fontsize{20}{23}\selectfont\bfseries #4\par}
  \end{minipage}\hfill
  \begin{minipage}[c]{2.6cm}\raggedleft\includegraphics[width=2.6cm]{\fm@logo}\end{minipage}\par
  \vspace{2pt}}

% ---------- barres de secció del color del curs ----------
\renewcommand{\seccio}[2]{\par\needspace{\fmneed\baselineskip}\bigskip
  \begin{tcolorbox}[colback=curs,colframe=curs,coltext=curstext,arc=2mm,boxrule=0pt,left=3mm,right=3mm,top=1.2mm,bottom=1.2mm]
  {\large\bfseries #1\quad #2}
  \end{tcolorbox}\vspace{-2pt}}

% ---------- estrelles de dificultat del color del curs (a Batxillerat no se'n posen) ----------
\@ifundefined{estrelles}{}{%
  \renewcommand{\estrelles}[1]{{\color{curs}\ifcase#1\or\faStar\faStar[regular]\faStar[regular]\or\faStar\faStar\faStar[regular]\or\faStar\faStar\faStar\fi}}}

% ---------- capçalera i peu de les pàgines ----------
% primera pàgina: logos institucionals al peu
\fancypagestyle{primera}{\fancyhf{}\renewcommand{\headrulewidth}{0pt}
  \fancyfoot[L]{\raisebox{-0.32\height}{\includegraphics[height=0.78cm]{logo-gva.png}}\hspace{5mm}%
    \raisebox{-0.32\height}{\includegraphics[height=0.78cm]{logo-vermellar.png}}}
  \fancyfoot[R]{\scriptsize\color{muted}\href{https://creativecommons.org/licenses/by-nc-sa/4.0/deed.ca}{\ccbyncsa}\ \ CC BY-NC-SA 4.0 · fentmates.org\qquad\small\thepage}}
% la resta: el símbol a l'esquerra i el curs a la dreta, amb el títol al centre com abans
\fancyhead[L]{\raisebox{-0.3\height}{\includegraphics[height=0.62cm]{\fm@marca}}}
\fancyhead[R]{\small\bfseries\color{cursfosc}\fm@cursnom}
\renewcommand{\headrule}{\color{curs!45}\hrule height \headrulewidth}
\makeatother

\fmcurs{1batx-ccss}

\begin{document}
\thispagestyle{primera}

\capcalera{Matemáticas CCSS}{1\textsuperscript{o} Bach.}{Ficha de ejercicios · Unidad 1}{\raggedright Números reales}
\vspace{10pt}
Nombre y apellidos: \rule{6.9cm}{0.4pt}\hfill Grupo: \rule{1.6cm}{0.4pt}\hfill Fecha: \rule{2.2cm}{0.4pt}
\vspace{8pt}

\begin{tcolorbox}[colback=soft,colframe=soft,arc=2mm,left=3mm,right=3mm,top=1.5mm,bottom=1.5mm]
\small Resuelve los ejercicios en el cuaderno mostrando todos los pasos. Da los resultados de manera exacta (fracciones y raíces simplificadas, sin raíces en el denominador) y, si hace falta, también aproximados. Escribe los conjuntos como intervalos y represéntalos en la recta real.
\end{tcolorbox}

% =====================================================
\seccio{1}{Números reales}

\exf{1} Indica el menor de los conjuntos numéricos ($\mathbb N$, $\mathbb Z$, $\mathbb Q$ o $\mathbb I$) a los que pertenecen los números siguientes:\begin{tasks}(5)
  \task $\sqrt5$
  \task $2$
  \task $-3$
  \task $24/6$
  \task $2,45$
  \task $-1,33$
  \task $2,3\widehat{45}$
  \task $\pi$
  \task $5/12$
  \task $\sqrt{\frac94}$
\end{tasks}

\ex{1} Razona si las afirmaciones siguientes son verdaderas o falsas:\begin{tasks}(1)
  \task Todo número entero es racional.
  \task Todo número decimal se puede expresar en forma de fracción.
  \task Todo número natural se puede escribir como una fracción.
  \task Si se suman dos números racionales, el resultado es siempre un número racional.
  \task Si se multiplican dos números irracionales, el resultado es siempre irracional.
\end{tasks}

\ex{1} Halla la fracción generatriz y simplifícala:\begin{tasks}(3)
  \task $0,1\widehat6$
  \task $2,\widehat{09}$
  \task $-1,2\widehat{45}$
  \task $3,\widehat{12}$
  \task $0,4\widehat5$
  \task $2,45$
\end{tasks}\recorda{Fracción generatriz: en el numerador, el número sin coma (hasta el final del primer periodo) menos la parte no periódica; en el denominador, tantos 9 como cifras tiene el periodo y tantos 0 como cifras tiene el anteperiodo.}

\ex{2} Realiza las operaciones siguientes escribiendo previamente los números en forma de fracción:\begin{tasks}(2)
  \task $2,\widehat7+3,\widehat4-1,\widehat1$
  \task $1,\widehat8+0,\widehat1$
  \task $0,\widehat3\cdot1,\widehat5-0,1\widehat6$
  \task $\left(1,\widehat2+0,5\right)\cdot0,\widehat{18}$
\end{tasks}

\ex{2} Clasifica cada número en el menor conjunto numérico. ¡Opera primero!\begin{tasks}(3)
  \task $\sqrt2\cdot\sqrt8$
  \task $\sqrt2+\sqrt8$
  \task $\pi-\pi$
  \task $0,1010010001\ldots$
  \task $-\sqrt{\frac94}$
  \task $\sqrt{-4}$
\end{tasks}


% =====================================================
\seccio{2}{Intervalos y semirrectas}

\exf{1} Completa la tabla siguiente:

\begin{center}\renewcommand{\arraystretch}{1.5}\begin{tabular}{|C{4.6cm}|C{3.2cm}|C{6cm}|}\hline\rowcolor{soft}\textbf{Desigualdad} & \textbf{Intervalo} & \textbf{Representación}\\\hline$\{x\in\mathbb R\ /\ x<-3\}$ & & \\\hline & $x\in[-2,\, 4)$ & \\\hline $\{x\in\mathbb R\ /\ -5<x\le-3\}$ & & \\\hline & $x\in(4,\, +\infty)$ & \\\hline\end{tabular}\end{center}

\ex{1} Escribe como intervalo o semirrecta y representa los conjuntos numéricos siguientes:\begin{tasks}(2)
  \task Números mayores que 3.
  \task $\{x\in\mathbb R/2\le x\le5\}$
  \task $\{x\in\mathbb R/3<x\le7\}$
  \task Números menores que 1 excluyendo el 0.
  \task Números reales positivos.
  \task $\{x\in\mathbb R/-2\le x<5\}$
  \task Números no negativos.
  \task $\{x\in\mathbb R/x\le-1\}$
\end{tasks}

\ex{1} Expresa como desigualdad, como intervalo y representa:\begin{tasks}(2)
  \task $x$ es menor que $-5$.
  \task 3 es menor o igual que $x$.
  \task $x$ está comprendido entre $-5$ y 1.
  \task $x$ está entre $-2$ y 0, ambos incluidos.
  \task $x$ es mayor o igual que $-1$ y menor que 6.
  \task $x$ no es mayor que 2.
\end{tasks}

\ex{1} Escribe la desigualdad que verifica todo número $x$ que pertenece a estos intervalos y semirrectas, y represéntalos:\begin{tasks}(4)
  \task $[-2,\, 7]$
  \task $[13,\, +\infty)$
  \task $(-\infty,\, 0)$
  \task $(-3,\, 0]$
  \task $\left[\frac32,\, 6\right)$
  \task $(-\infty,\, +\infty)$
  \task $(1,\, 6)$
  \task $(0,\, 4]$
  \task $(-\infty,\, 2)$
  \task $(-\infty,\, {-}5]$
  \task $(-1,\, +\infty)$
  \task $\left(2,\, \frac72\right)$
\end{tasks}

\ex{1} ¿Qué números enteros pertenecen a cada uno de los conjuntos siguientes?\begin{tasks}(4)
  \task $[-2,\, 3)$
  \task $\left(-1,\, \frac52\right]$
  \task $(-4,\, {-}1)$
  \task $(3,\, 4)$
\end{tasks}


% =====================================================
\seccio{3}{Valor absoluto}

\exf{1} Calcula:\begin{tasks}(4)
  \task $|-7|$
  \task $|3-8|$
  \task $-|-5|$
  \task $|-2|+|5|$
  \task $|(-3)\cdot4|$
  \task $\big||-5|-8\big|$
  \task $|\pi-4|$
  \task $|1-\sqrt3|$
\end{tasks}

\ex{1} Calcula la distancia entre los pares de números siguientes:\begin{tasks}(4)
  \task $-4$ y $6$
  \task $2$ y $7$
  \task $-3$ y $-9$
  \task $-\frac12$ y $\frac32$
\end{tasks}

\ex{2} Resuelve las ecuaciones siguientes:\begin{tasks}(3)
  \task $|x|=5$
  \task $|x-3|=2$
  \task $|2x+1|=7$
  \task $|x+4|=-1$
  \task $|3x-6|=0$
  \task $|5-x|=1$
\end{tasks}\recorda{Si $r>0$, \ $|A|=r\iff A=r$ o $A=-r$.}

\ex{2} Averigua para qué valores de $x$ se cumplen las desigualdades siguientes. Escribe la solución como intervalo y represéntala:\begin{tasks}(3)
  \task $|x|<3$
  \task $|x|\ge3$
  \task $|x-2|\le3$
  \task $|x-4|>2$
  \task $|2x-1|>1$
  \task $|-5x+3|<7$
  \task $|x+1|\le\frac12$
  \task $|-5x+25|\le-2$
  \task $\left|x+\frac34\right|>-1$
\end{tasks}\recorda{Si $r>0$: \ $|A|<r\iff -r<A<r$ \quad y \quad $|A|>r\iff A<-r$ o $A>r$. Al multiplicar o dividir por un número negativo, la desigualdad cambia de sentido.}

\ex{2} Expresa cada conjunto con una desigualdad con valor absoluto (busca el centro y el radio):\begin{tasks}(4)
  \task $(1,\, 7)$
  \task $[-5,\, 1]$
  \task $(-\infty,\, {-}1)\cup(5,\, +\infty)$
  \task $\left[-\frac12,\, \frac12\right]$
\end{tasks}


% =====================================================
\seccio{4}{Unión e intersección de intervalos}

\exf{2} Dados los intervalos $A=[-3,\, 2]$, $B=[0,\, 5]$, $C=[2,\, +\infty)$, $D=(0,\, 10)$ y $E=\left(2,\, \frac72\right)$, calcula:\begin{tasks}(3)
  \task $A\cap B$
  \task $A\cap D$
  \task $A\cup B$
  \task $(A\cup B)\cap C$
  \task $A\cap C$
  \task $B\cup C$
  \task $C\cap D$
  \task $E\cap B$
  \task $A\cup E$
\end{tasks}\recorda{Dibuja los intervalos uno encima del otro en la misma recta. La unión ($\cup$) son los números que están en alguno de los dos; la intersección ($\cap$), los que están en los dos a la vez.}

\ex{2} Dados $A=(-\infty,\, {-}4]$, $B=[2,\, 5)$, $C=[-6,\, 1]$, $D=[0,\, +\infty)$ y $E=(-2,\, {-}1)$, calcula:\begin{tasks}(3)
  \task $A\cup B$
  \task $B\cap C$
  \task $D\cap C$
  \task $A\cup(B\cap D)$
  \task $B\cap(A\cup E)$
  \task $B\cap D\cap E$
  \task $(B\cup C)\cap A$
  \task $A\cup C$
  \task $C\cap E$
  \task $A\cap D$
\end{tasks}

\ex{1} Escribe en forma de intervalo los números que verifican:\begin{tasks}(3)
  \task $x<3$ y $x\ge5$
  \task $x>0$ y $x<4$
  \task $x\le-1$ y $x>1$
  \task $x<3$ y $x\ge-2$
  \task $x<-1$ o $x\ge2$
  \task $x>0$ o $x<4$
\end{tasks}

\ex{2} Sea $A$ el conjunto de los números que cumplen $|x-1|<2$ y $B=[2,\, +\infty)$. Escribe $A$ como intervalo y calcula $A\cup B$ y $A\cap B$. Represéntalos.


% =====================================================
\seccio{5}{Radicales}

\exf{1} Calcula el valor de los radicales siguientes, si existen:\begin{tasks}(4)
  \task $\sqrt{16}$
  \task $\sqrt[3]{-8}$
  \task $\sqrt[5]{32}$
  \task $\sqrt[3]{\frac{27}{8}}$
  \task $\sqrt[3]{0,064}$
  \task $\sqrt[4]{81}$
  \task $\sqrt[4]{-16}$
  \task $\sqrt[5]{-\frac{1}{32}}$
\end{tasks}

\ex{1} Escribe en forma de potencia los radicales de la primera fila y en forma de radical las potencias de la segunda:\begin{tasks}(5)
  \task $\sqrt5$
  \task $\sqrt[3]{x^2}$
  \task $\sqrt[5]{a^3}$
  \task $\dfrac{1}{\sqrt2}$
  \task $\dfrac{1}{\sqrt[3]{a^2}}$
  \task $3^{\frac25}$
  \task $a^{\frac34}$
  \task $5^{-\frac12}$
  \task $x^{1,5}$
  \task $2^{\frac53}$
\end{tasks}\recorda{$\sqrt[n]{a^m}=a^{\frac mn}$}

\ex{1} Calcula sin calculadora pasando a exponentes fraccionarios:\begin{tasks}(3)
  \task $8^{\frac23}$
  \task $81^{0,75}$
  \task $\left(\dfrac{1}{32}\right)^{-\frac25}$
  \task $16^{\frac34}$
  \task $27^{-\frac13}$
  \task $25^{\frac32}$
\end{tasks}

\ex{1} Simplifica los radicales siguientes:\begin{tasks}(4)
  \task $\sqrt[6]{2^4}$
  \task $\sqrt[9]{a^6}$
  \task $\sqrt[15]{8a^3}$
  \task $\sqrt[8]{16a^4}$
  \task $\sqrt[6]{0,027}$
  \task $\sqrt[12]{x^8y^4}$
  \task $\sqrt[4]{25}$
  \task $\sqrt[6]{64}$
\end{tasks}

\ex{2} Reduce los radicales de cada apartado al mínimo índice común:\begin{tasks}(2)
  \task $\sqrt2,\ \sqrt[3]3$
  \task $\sqrt[3]{n^2},\ \sqrt[4]{m^3},\ \sqrt[6]{n^5}$
  \task $\sqrt a,\ \sqrt[3]{a^2},\ \sqrt[4]{a^3}$
  \task $\sqrt[3]{5xy^2},\ \sqrt{6x^3z}$
\end{tasks}

\ex{2} Ordena de menor a mayor reduciendo a índice común:\begin{tasks}(3)
  \task $\sqrt3,\ \sqrt[3]5,\ \sqrt[4]{10}$
  \task $\sqrt2,\ \sqrt[3]3,\ \sqrt[6]6$
  \task $\sqrt[3]2,\ \sqrt[4]3$
\end{tasks}

\ex{2} Extrae todos los factores posibles de cada uno de los radicales siguientes:\begin{tasks}(5)
  \task $\sqrt8$
  \task $\sqrt{12}$
  \task $\sqrt[3]{16}$
  \task $\sqrt[3]{54}$
  \task $\sqrt[5]{64}$
  \task $\sqrt{\frac{27}{4}}$
  \task $3\sqrt{8a^3}$
  \task $\sqrt[3]{54a^4}$
  \task $\sqrt{72x^3y^4}$
  \task $\sqrt[4]{32x^6}$
\end{tasks}

\ex{2} Introduce el coeficiente dentro del radical y simplifica:\begin{tasks}(3)
  \task $7\sqrt a$
  \task $2a\sqrt{3a}$
  \task $x\sqrt{\frac1x}$
  \task $\frac32\sqrt{\frac23}$
  \task $2\sqrt[3]3$
  \task $3\sqrt[3]2$
  \task $\frac23\sqrt[3]{\frac{81}{4}}$
  \task $a^2\sqrt a$
  \task $3+\sqrt7$
\end{tasks}

\ex{2} Realiza las sumas y restas siguientes:\begin{tasks}(2)
  \task $\sqrt8+\sqrt{18}+\sqrt{32}+\sqrt{50}$
  \task $\sqrt{28}+\sqrt{63}-\sqrt{175}$
  \task $3\sqrt2-3\sqrt8+3\sqrt{18}$
  \task $2a\sqrt3-\sqrt{27a^2}+a\sqrt{12}$
  \task $5\sqrt{20}-2\sqrt{45}+\sqrt5$
  \task $\sqrt{18}+5\sqrt{12}-2\sqrt2-\sqrt{243}$
  \task $\sqrt{\frac13}+\sqrt{27}$
  \task $2\sqrt[3]{16}-\sqrt[3]{54}+\sqrt[3]{250}$
\end{tasks}\recorda{Solo se pueden sumar radicales semejantes (mismo índice y mismo radicando). Antes, extrae sus factores.}

\ex{2} Calcula y simplifica (reduce a índice común cuando haga falta):\begin{tasks}(4)
  \task $\sqrt5\cdot\sqrt6$
  \task $\sqrt{12}\cdot\sqrt3$
  \task $\sqrt[3]4\cdot\sqrt[3]{16}$
  \task $\sqrt5\cdot\sqrt[3]6$
  \task $\sqrt{20}:\sqrt5$
  \task $\sqrt[4]{9x^5y^3}:\sqrt[4]{3x^3y}$
  \task $\dfrac{\sqrt2}{\sqrt[3]2}$
  \task $\sqrt a\cdot\sqrt[3]a\cdot\sqrt[4]a$
\end{tasks}

\ex{2} Calcula las potencias y escribe bajo un solo radical:\begin{tasks}(3)
  \task $(\sqrt3)^4$
  \task $(\sqrt[3]2)^6$
  \task $(\sqrt2)^5$
  \task $\left(\sqrt[4]{a^3}\right)^2$
  \task $(2\sqrt3)^2$
  \task $\sqrt[3]{(27x^3y^6)^5}$
  \task $\sqrt{\sqrt8}$
  \task $\sqrt[3]{\sqrt{32}}$
  \task $\sqrt{2\sqrt2}$
\end{tasks}

\ex{2} Opera y simplifica:\begin{tasks}(3)
  \task $(\sqrt3+\sqrt2)^2$
  \task $(2\sqrt3-1)(2\sqrt3+1)$
  \task $(\sqrt5-2)^2$
  \task $\sqrt2\left(\sqrt8-\sqrt{18}\right)$
  \task $(1+\sqrt2)^2-\sqrt8$
\end{tasks}


% =====================================================
\seccio{6}{Racionalización}

\exf{2} Racionaliza el denominador de las expresiones siguientes:\begin{tasks}(5)
  \task $\dfrac{5}{\sqrt5}$
  \task $\dfrac{2\sqrt6}{\sqrt2}$
  \task $\dfrac{3}{\sqrt5}$
  \task $\sqrt{\dfrac52}$
  \task $\dfrac{3}{2\sqrt[3]4}$
  \task $\dfrac{6}{\sqrt[3]3}$
  \task $\dfrac{3y}{x\sqrt y}$
  \task $\dfrac{2\sqrt{12}}{5\sqrt3}$
  \task $\dfrac{10}{\sqrt[4]8}$
\end{tasks}\recorda{Si en el denominador hay una sola raíz, multiplicamos arriba y abajo por lo que haga falta para que la raíz desaparezca: $\dfrac{1}{\sqrt[n]{a^m}}=\dfrac{\sqrt[n]{a^{n-m}}}{a}$.}

\ex{2} Racionaliza el denominador de las expresiones siguientes:\begin{tasks}(3)
  \task $\dfrac{1}{\sqrt2+\sqrt3}$
  \task $\dfrac{2}{\sqrt2-1}$
  \task $\dfrac{\sqrt2}{3-\sqrt2}$
  \task $\dfrac{1-\sqrt2}{1+\sqrt2}$
  \task $\dfrac{4}{\sqrt5-1}$
  \task $\dfrac{5}{2(\sqrt3-\sqrt2)}$
  \task $\dfrac{2+\sqrt2}{2-\sqrt2}$
  \task $\dfrac{\sqrt3}{\sqrt3+1}$
  \task $\dfrac{1}{\sqrt x-\sqrt y}$
\end{tasks}\recorda{Si en el denominador hay una suma o una resta con raíces, multiplicamos por el conjugado: $(a+b)(a-b)=a^2-b^2$.}

\ex{2} Racionaliza, opera y simplifica:\begin{tasks}(3)
  \task $\dfrac{\sqrt3+1}{\sqrt3-1}-\dfrac{\sqrt3-1}{\sqrt3+1}$
  \task $\dfrac{1}{\sqrt2}+\dfrac{1}{\sqrt2+1}$
  \task $\dfrac{2}{\sqrt3}-\dfrac{1}{2+\sqrt3}$
\end{tasks}


% =====================================================
\seccio{7}{Logaritmos}

\exf{1} Aplicando la definición, calcula el valor de los logaritmos siguientes:\begin{tasks}(4)
  \task $\log_2 8$
  \task $\log_3 81$
  \task $\log 1000$
  \task $\log 0,0001$
  \task $\ln e^{5}$
  \task $\log_4 16$
  \task $\log_4 0,25$
  \task $\log_5 1$
\end{tasks}\recorda{$\log_a b=c\iff a^c=b$ \ ($a>0$, $a\neq1$, $b>0$). $\log$ es el logaritmo en base 10 y $\ln$, en base $e$.}

\ex{1} Aplicando la definición, calcula el valor de los logaritmos siguientes:\begin{tasks}(4)
  \task $\log_3\sqrt{27}$
  \task $\log_7\frac{1}{49}$
  \task $\log_3 3\sqrt3$
  \task $\log_3 0,\widehat3$
  \task $\log_5\sqrt[3]{25}$
  \task $\log_9\sqrt[3]3$
  \task $\log_{\frac12}8$
  \task $\log_8 0,125$
\end{tasks}

\ex{2} Calcula, si es posible, el valor de $x$ en cada una de las expresiones siguientes:\begin{tasks}(4)
  \task $\log_x 8=3$
  \task $\log_x 8=-3$
  \task $\log_5 x=-1$
  \task $\log_2 x=5$
  \task $\log_x 49=2$
  \task $\log_3(-81)=x$
  \task $\log_x\sqrt2=0$
  \task $\log x=-2$
\end{tasks}\recorda{$\log_a b$ solo existe si $a>0$, $a\neq1$ y $b>0$.}

\ex{1} Con la ayuda de la calculadora y el cambio de base, obtén aproximaciones hasta las milésimas:\begin{tasks}(3)
  \task $\log_3 20$
  \task $\log_{0,5}60$
  \task $\log_{\sqrt2}3$
  \task $\log_7 100$
  \task $\log_2 0,3$
  \task $\log_5 2$
\end{tasks}\recorda{Cambio de base: $\log_a x=\dfrac{\log x}{\log a}=\dfrac{\ln x}{\ln a}$.}


% =====================================================
\seccio{8}{Propiedades de los logaritmos}

\exf{1} Halla el resultado de las expresiones siguientes mediante las propiedades de los logaritmos:\begin{tasks}(2)
  \task $2\log_4 16+\log_2 32-3\log_7 49$
  \task $\log_2 8+\log_3 27+3\log_5 125$
  \task $\log_5 625-\log_9 81+\log_8 64$
  \task $\log_6 4+\log_6 9$
  \task $\log_2 40-\log_2 5$
  \task $\log 20+\log 5$
  \task $\log_3 54-\log_3 2$
  \task $2\log 5+\log 4$
\end{tasks}\recorda{$\log_a(p\cdot q)=\log_a p+\log_a q$, \ $\log_a\frac pq=\log_a p-\log_a q$, \ $\log_a p^n=n\log_a p$.}

\ex{2} Considerando $\log2=0,301$ y $\log3=0,477$, calcula sin calculadora:\begin{tasks}(5)
  \task $\log4$
  \task $\log6$
  \task $\log5$
  \task $\log12$
  \task $\log1,5$
  \task $\log0,02$
  \task $\log3,6$
  \task $\log\sqrt3$
  \task $\log0,\widehat3$
  \task $\log45$
\end{tasks}

\ex{2} Si $\log_3 2=0,631$ y $\log_3 5=1,465$, calcula $\log_3 150$ y $\log_3\frac{10}{9}$.

\ex{2} Toma logaritmos decimales y desarrolla las expresiones:\begin{tasks}(2)
  \task $P=10x^3y$
  \task $Q=\dfrac{100x^2}{x+y}$
  \task $R=\dfrac{x^2\sqrt y}{z^3}$
  \task $S=\sqrt[3]{a^2b}$
\end{tasks}

\ex{2} Expresa el valor de $E$ en cada caso sin que aparezcan logaritmos:\begin{tasks}(1)
  \task $\log E=2-2\log x+\log y$
  \task $\log E=3\log2-4\log x+3\log y$
  \task $\log E=\log(x-2y)+\log(x+y)$
  \task $\log E=\frac12\log x+\log3$
  \task $\log E=1+\log a-2\log b$
\end{tasks}


% =====================================================
\seccio{9}{Problemas}

\exf{2} La temperatura $T$ de una nevera tiene que ser de $4\,^\circ$C con un margen de error de $1,5\,^\circ$C. Escribe esta condición con un valor absoluto y como intervalo. ¿Es aceptable una lectura de $5,7\,^\circ$C?

\ex{2} Un cuadrado tiene un área de $50\ \text{cm}^2$. Calcula su lado y su diagonal en forma de radical simplificado.

\ex{2} Ingresamos 2\,000 € en una cuenta al 3\,\% de interés compuesto anual. ¿Cuántos años tienen que pasar, como mínimo, para que haya 3\,000 €?

\ex{2} Un coche cuesta 20\,000 € y pierde un 15\,\% de su valor cada año. ¿Cuántos años tienen que pasar para que valga menos de 8\,000 €?

\ex{2} Una ciudad tiene 50\,000 habitantes y la población crece un 2\,\% cada año. ¿Cuántos años tardará en llegar a 60\,000 habitantes?

\ex{2} Una bacteria se duplica cada hora si tiene las condiciones óptimas. ¿Cuántas horas tendrán que pasar para que una bacteria inicial cree una colonia de mil millones de bacterias?

\ex{1} Encuentra el error en cada uno de estos cálculos y corrígelo:\begin{tasks}(2)
  \task $\sqrt{9+16}=3+4=7$
  \task $-3^2=9$
  \task $\log(5+5)=\log5+\log5$
  \task $2^{-3}=-8$
  \task $\sqrt2+\sqrt3=\sqrt5$
  \task $[1,\, 3)\cap(3,\, 5]=\{3\}$
\end{tasks}

\ex{3} Sin calculadora, ordena de menor a mayor $\log_2 10$, \ $3$ \ y \ $\log_3 10$. Justifícalo.

\vfill
\begin{center}\color{muted}\small\faCheckCircle[regular]\ \ ¡Ánimo! Comprueba tus resultados con las herramientas interactivas de los apuntes de la unidad.\end{center}

\end{document}
